\documentclass[a4paper,11pt]{article}
\usepackage[utf8]{inputenc}
\usepackage[T1]{fontenc}
\usepackage[english]{babel}
\usepackage{booktabs}
\usepackage{tabularx}
\usepackage{listings}
\usepackage{hyperref}
\usepackage{geometry}
\geometry{margin=2.5cm}

\title{Memory Annex --- Extracted Sections}
\author{Eloi Egea Rada}
\date{TFG 2025--2026}

\begin{document}
\maketitle
\tableofcontents
\newpage

% ============================================================
% HOW TO USE THIS FILE
% Each \section corresponds to a chapter of the main memory.
% Removed blocks are enclosed in \subsection with a note
% explaining why they were moved here.
% ============================================================

% ============================================================
\section{Chapter 1 --- Introduction}
% ============================================================

\subsection{Document Structure (removed from Ch.1)}

% Reason: Standard boilerplate that adds length without contributing
% to the academic argument. The chapter ordering is self-evident
% from the table of contents. Removed 2026-04-27.

The remainder of this document is organised as follows.
Chapter~2 reviews the state of the art in cybersecurity education platforms and
lab environments, situating this project within the existing landscape.
Chapter~3 presents the feasibility study, covering the project timeline, budget,
technical viability, and risk assessment.
Chapter~4 defines the project objectives and measurable KPIs.
Chapter~5 describes the methodology, project phases, and research methods.
Chapter~6 specifies the functional and technical requirements of the lab package.
Chapter~7 documents the development process, including platform setup, infrastructure
decisions, and a detailed description of each implemented lab.

% ============================================================
\section{Chapter 2 --- State of the Art}
% ============================================================

\subsection{Removed from Ch.2: Platform Descriptions (7 subsections)}

% Reason: The individual platform descriptions are subsumed by the
% comparative table (Table 2.1). Removing them reduces redundancy
% and keeps the chapter focused on analysis rather than description.
% Removed 2026-04-27.

\subsubsection{HackTheBox}
HackTheBox \cite{hackthebox} is one of the most widely recognised online platforms for
penetration testing practice, offering a large library of intentionally vulnerable machines and
Capture-the-Flag challenges categorised by difficulty. The platform benefits from a strong
community and industry recognition, making it a common supplementary resource for
professional certifications. However, its challenges are isolated: there is no network topology
connecting systems, no formal curriculum mapping to specific degree programmes, and the
platform cannot be deployed on institutional infrastructure, which limits customisation and
assessment integration.

\subsubsection{TryHackMe}
TryHackMe \cite{tryhackme} adopts a gamified, beginner-friendly approach, combining structured
learning paths with guided interactive labs accessible entirely from a browser. Its low barrier to
entry and progressive structure make it effective for initial skill development. However, the
platform is externally managed and subscription-dependent, scenarios are simplified to ensure
completion within short timeframes, and it does not support local institutional deployment or
integration with academic learning management systems.

\subsubsection{Cisco Networking Academy}
Cisco Networking Academy \cite{cisco_netacad} provides institution-backed curriculum aligned with
professional certifications (CCNA, CyberOps) and uses Packet Tracer for network simulation.
Its formal structure and vendor support are strengths for foundational networking education.
However, the programme is primarily defensive and networking-oriented, Packet Tracer does not
support full operating system simulation or security tool integration, and the curriculum is
tightly coupled to Cisco technologies.

\subsubsection{SEED Labs}
SEED Labs \cite{seedlabs} is an open-source, research-backed collection of Docker-based security
exercises covering specific concepts such as buffer overflows, SQL injection, and cryptographic
protocols. The academic rigour and local deployability are significant advantages. Nevertheless,
each lab is conceptually isolated, there is no structured penetration testing workflow spanning
multiple systems, and significant instructor customisation is required to integrate SEED exercises
into a complete curriculum.

\subsubsection{Game of Active Directory (GOAD)}
GOAD \cite{goad} provides an Infrastructure-as-Code vulnerable Active Directory environment
designed for realistic Windows domain penetration testing. Its use of Terraform and Vagrant enables
reproducible deployment and the scenarios reflect real corporate environments. The platform is,
however, narrowly focused on Windows AD: it does not cover reconnaissance, web application
vulnerabilities, or network traffic analysis, and it lacks formal instructor resources or
curriculum structure.

\subsubsection{VulnHub}
VulnHub \cite{vulnhub} is a community-driven repository of downloadable vulnerable virtual machines
for local deployment. Its open, self-directed nature suits independent learners well, but the
absence of structured progression, inconsistent quality across community-contributed machines,
and lack of automation for deployment and reset make institutional integration impractical
without significant instructor effort.

\subsubsection{Hack4u Academy}
Hack4u Academy \cite{hack4u} stands out for its emphasis on penetration testing methodology over
tool memorisation, offering structured video-based courses with integrated practical exercises.
Its curriculum is well-designed and the community active. The platform is, however,
subscription-based and externally hosted, making local institutional deployment and automated
large-scale assessment impossible.

% ============================================================
\section{Chapter 3 --- Objectives and Scope}
% ============================================================

\subsection{Removed from Ch.3: OE Deliverables blocks}

% Reason: Deliverables are already captured in the project repository
% and src/materials/exercises/ structure. Listing them in the memory
% duplicates information and inflates the objectives section.
% Removed 2026-04-27.

\subsubsection{OE1 Deliverables}
\begin{itemize}
\item 4 .unl topology files (.unl), one per lab
\item 5 core virtual machine images (VyOS, pfSense~CE, Ubuntu Server~24.04,
      Ubuntu Desktop~24.04, Parrot~OS~Security), sized and optimised for
      concurrent EVE-NG deployment on constrained HomeLab hardware
\item Technical configuration documentation for each lab
\end{itemize}

\subsubsection{OE2 Deliverables}
\begin{itemize}
\item Bash/Python utility scripts for common tasks (e.g.\ bridge configuration,
      ISO upload, node connectivity checks)
\item Scripts documented with usage instructions in the repository
\end{itemize}

\subsubsection{OE3 Deliverables}
\begin{itemize}
\item Student lab guide PDF (enunciado) per lab --- task description, objectives,
      hints, and tools reference
\item Instructor resolution guide PDF (resolució) per lab --- full walkthrough,
      flag values, rubric (100~pts), and common errors
\item Technical configuration reference per lab in the repository
\end{itemize}

\subsubsection{OE4 Deliverables}
\begin{itemize}
\item Pilot session report with timing and usability observations
\item Post-session survey results (satisfaction, difficulty, clarity)
\item List of improvements implemented based on pilot feedback
\end{itemize}

\subsection{Removed from Ch.3: Secondary and Potential Audiences}

% Reason: Too generic for the scope of the TFG. The primary audience
% (professor + GEISI students) is sufficient for the objectives chapter.
% Removed 2026-04-27.

\subsubsection{Secondary End Users}
\begin{itemize}
\item \textbf{Cybersecurity Faculty}: Other teachers interested in reusing the material
\item \textbf{Postgraduate Students}: Possible extensions of the material for advanced levels
\end{itemize}

\subsubsection{Potential Audience --- Internal Scope (Tecnocampus)}
\begin{itemize}
\item Other courses related to cybersecurity
\item Research projects in computer security
\item Continuing education and professional certifications
\end{itemize}

\subsubsection{Potential Audience --- External Scope}
\begin{itemize}
\item Educational institutions with training in cybersecurity
\item Specialized vocational training centers
\item Companies with internal training programs in security
\end{itemize}

\subsection{Removed from Ch.3: Three separate KPI tables (merged into one)}

% Reason: Three tables with identical structure were redundant.
% Merged into a single table with a Category column.
% Removed 2026-04-27.

% (original tables: tab:kpi_technical, tab:kpi_educational, tab:kpi_quality)

% ============================================================
\section{Chapter 4 --- Methodology}
% ============================================================

\subsection{Removed from Ch.4: Detailed phase activity lists and deliverables}

% Reason: The numbered sub-activity lists within each phase were verbose
% and duplicated content already in Ch.3 (OE deliverables). Condensed
% each phase to an objective + summary sentence + duration/status.
% Removed 2026-04-27.

\subsubsection{Phase 1 detailed activities (8 numbered items)}
\begin{enumerate}
\item SMART Objectives Definition (4h): Specification of MO, 4 OEs with KPIs, GEISI alignment
\item State of the Art Analysis (15h): 7 platforms reviewed, gap analysis, comparative matrix
\item Functional and Non-Functional Requirements (10h): RF1-RF3, NFR1-NFR3
\item Pentesting Methodology and Validation Framework (8h): PTES adaptation, CEH principles
\item Design Decisions Justification (5h): EVE-NG rationale, tech stack
\item Risk Analysis and Mitigation Planning (15h): 12+ risks identified
\item Resource Inventory and Allocation (5h): hardware, software, human resources, budget
\item Gantt Chart and Critical Path Analysis (20h): 37 tasks, 13 critical path tasks
\end{enumerate}

\subsubsection{Phase 2 detailed activities (6 numbered items)}
\begin{enumerate}
\item Lab 1 Reconnaissance (14h research + 20h dev): PTES, Nmap, multi-tier topology
\item Lab 2 Web Vulnerabilities (14h research + 25h dev): OWASP Top 10, Burp, ZAP
\item Lab 3 Network Analysis (14h research + 25h dev): NIST, Wireshark, multi-VLAN
\item Pilot Validation (20h + 6h revision): 4--5 GEISI students, surveys, adjustments
\item Topology Adjustments (30h): performance tuning, VM optimisation
\item Teaching Material (40h): enunciado + resolució PDFs per lab
\end{enumerate}

\subsubsection{Phase 3 detailed activities (5 numbered items)}
\begin{enumerate}
\item Lab 4 Privilege Escalation (16h research + 30h dev): CEH, Metasploit, hardened systems
\item Comprehensive Penetration Testing (30h Round 1 + 25h Round 2): all 4 labs
\item Automation Scripts (30h): Bash/Python utilities
\item Final User Validation (18h): pilot group, surveys, final adjustments
\item Report Writing (90h): Phases I--III + revision + formatting
\end{enumerate}

\subsubsection{Phase Deliverables blocks}
% Each phase had a \textbf{Deliverables} itemize block --- removed since
% deliverables are captured in OE1-OE4 (Ch.3) and the project repository.

\subsection{Removed from Ch.4: Information Search Strategies detail}

% Reason: Condensed to 1 paragraph. The itemized database lists and
% tool doc references were generic boilerplate.

\begin{itemize}
  \item Academic databases: IEEE Xplore, ACM Digital Library, Scopus, Google Scholar
  \item Industry standards: PTES, OWASP, NIST SP800-115, CIS Controls, EC-Council CEH
  \item Tool documentation: EVE-NG, Parrot OS, Metasploit, OWASP ZAP, Burp Suite
  \item Comparative analysis: feature matrices, user reviews, educational effectiveness studies
\end{itemize}

Primary Research Methods:
\begin{itemize}
  \item User Research: informal pilot sessions, qualitative interviews, quantitative metrics
  \item Technical Experimentation: lab iteration, performance testing, security validation
\end{itemize}

\subsection{Removed from Ch.4: QA subsections and Documentation subsection}

% Reason: Generic content applicable to any software project; not specific
% enough to add value in a cybersecurity TFG context.

Quality Assurance covered: Technical Validation (automated + manual testing, performance
benchmarking, compatibility testing), Pedagogical Validation (GEISI outcome alignment,
rubric assessment, faculty feedback), and Usability Testing (UI/doc clarity, accessibility,
user satisfaction scoring).

Documentation and Version Control: GitHub repository, LaTeX for formal docs, versioning
system, change logs, deployment guides and runbooks.

% ============================================================
\section{Chapter 5 --- Functional and Non-Functional Requirements}
% ============================================================

\subsection{Removed from Ch.5: Generic RF sub-requirements}

% Reason: RF1.2-RF1.4, RF2.5, RF3.3-RF3.4 are generic software
% engineering specifications not specific to this cybersecurity project.
% Removed 2026-04-27.

RF1.2 (5-8 VMs per topology, realistic network arch, VLANs),
RF1.3 (EVE-NG compatibility, KVM, Ubuntu 20.04+, nested virt),
RF1.4 (topology documentation: diagrams, IP scheme, port mappings, attack surface),
RF2.5 (vulnerability implementation standards: fixable, reset procedure, no unrecoverable systems),
RF3.3 (support materials: student manuals, teacher guides, troubleshooting),
RF3.4 (validation materials: automated scripts, manual procedures, success indicators).

\subsection{Removed from Ch.5: Generic TR sub-requirements}

% Reason: TR1.2, TR1.4, TR2.4, TR3.3 are generic infrastructure/compliance
% items not specific to the project.

TR1.2 (Hypervisor: KVM, Ubuntu Server 20.04/22.04, QEMU, nested virt),
TR1.4 (Network config: L2 switching, VLAN, IPv4/IPv6),
TR2.4 (Tool compliance: free/educational licenses, open-source preference, legal guidelines),
TR3.3 (Integration: EVE-NG API, Docker for SYNAPSE, Git/GitHub).

\subsection{Removed from Ch.5: NFR1-NFR6 detailed subsections (replaced by table)}

% Reason: Six NFR subsections with 3 sub-requirements each totalling ~3 pages
% were replaced by a single compact table (Table 5.1). All key values preserved.

NFR1 Performance --- NFR1.1 Deployment (deploy ≤2min, boot ≤90s, reset ≤30s, 5 concurrent),
NFR1.2 Resource (≤5GB images, ≤6GB RAM/lab), NFR1.3 Scalability (3-4 concurrent HomeLab,
horizontal scaling, EVE-NG Pro for classroom).

NFR2 Compliance --- NFR2.1 NIST CSF (5 functions mapped to labs), NFR2.2 CIS Controls
(Controls 1-18, CIS benchmarks in rubrics), NFR2.3 OWASP (Top 10 full, testing guide,
dev practices), NFR2.4 PTES+CEH (phases mapped, ethical principles enforced).

NFR3 Reliability --- NFR3.1 System (MTBF≥100h, MTTR≤5min, 100\% reproducibility, 0
critical bugs), NFR3.2 Data Integrity (no data loss, snapshot/restore, config backup),
NFR3.3 Availability (maintenance weekends only, ≥99\% during instructional periods).

NFR4 Security --- NFR4.1 Sandbox (network isolation, no outbound internet, contained attacks),
NFR4.2 Access Control (RBAC, audit logs, authorized/unauthorized identification),
NFR4.3 Ethics (mandatory training, written agreements, monitoring/escalation).

NFR5 Usability --- NFR5.1 Ease of Use (5-10 min faculty deploy, intuitive UI, minimal
student prereqs), NFR5.2 Documentation (user manuals, quick-start; video tutorials and
FAQ not produced), NFR5.3 Accessibility (WCAG 2.1 AA, screen reader, keyboard nav).

NFR6 Sustainability --- NFR6.1 Maintainability (all open-source, no vendor lock-in,
GitHub source, backward compat), NFR6.2 Curriculum (GEISI outcomes mapping, flexible
customisation, multi-level courses), NFR6.3 Scalability (institutional infra, multi-faculty,
scheduling, reporting).

% ============================================================
\section{Chapter 6 --- Feasibility Study}
% ============================================================

\subsection{Removed from Ch.6: Gantt phase breakdown (task lists)}

% Reason: Redundant with Ch.4 (Methodology) phase descriptions.
% Replaced with a reference to Chapter 4. Removed 2026-04-27.

Phase 0 tasks: 0.1 HomeLab Init (64h), 0.2 VSCode Setup (18h),
0.3 Parrot OS Setup (15h), 0.4 As-Built Documentation (8h).

Phase 1 tasks: 1.0 Thesis Dev (22h), 1.1 SMART Objectives (4h),
1.2 State of Art (15h), 1.3 Requirements (10h), 1.4 Pentesting Method (8h),
1.5 Design Decisions (5h), 1.6 Gantt (20h), 1.7 Risk Analysis (15h),
1.8 Resource Inventory (5h), Milestone 1.10 Preliminary Project (40h).

Phase 2 tasks: 2.1-2.3 Lab research (14h each), 2.4 Write-ups (13h),
2.5 Topology Adj (30h), 2.6 Pre-validation (20h), 2.6.1 Revision (6h),
2.7 Hackathon (17h), 2.8 Interim Report (40h), 2.9 Servers/VMs setup (20h),
2.10 Backup (10h), Milestone 2.11 (50h).

Phase 3 tasks: 3.1 Lab4 Research (16h), 3.2 Lab3 Validation (12h),
3.3 Write-ups (19h), 3.4-3.5 Pentest Rounds (30h+25h), 3.6 Scripts (30h),
3.7 Final Validation (18h), 3.8 Polishing (25h), 3.9-3.11 Report Writing (90h),
3.12 Revision (15h), 3.13 Preparation (20h), Milestone 3.14 (5h).

\subsection{Removed from Ch.6: NIST 5-phase risk assessment detail}

% Reason: ~200 lines covering the 5 NIST SP 800-30 phases in full detail.
% Replaced by the compact risk matrix table (Table 6.4) + conclusion.
% Removed 2026-04-27.

Phase 1 Asset Identification: Hardware assets (HomeLab server, workstations,
network equipment, storage), Software assets (EVE-NG, Parrot OS, dev tools),
Data assets (thesis source, lab configs, learning materials), Infrastructure assets
(virtual network, VLANs, backups). Scope excludes organizational and regulatory risks.

Phase 2 Threat/Vulnerability Identification: Hardware threats (disk failure, power,
overheating, RAM exhaustion), Software threats (outdated images, config drift, EVE-NG
crashes), Knowledge gaps, Resource constraints, Integration threats (EVE-NG/host
compatibility), Data loss. Vulnerabilities: aging components, single-point-of-failure,
community edition limitations, inadequate backup redundancy.

Phase 3 Impact/Likelihood: Scales Low/Medium/High for both axes per NIST SP 800-30.

Phase 4 Prioritization: HIGH = EVE-NG Performance; MEDIUM = VM Incompatibility,
Knowledge Gaps, Config Documentation, Data Loss; LOW = Hardware Degradation.

Phase 5 Mitigation strategies in full detail (REDUCTION/PREVENTION strategies,
concrete actions, result probability reductions).

\subsection{Removed from Ch.6: Economic cost-benefit bullets}

% Reason: Generic bullet points subsumed by the cost comparison table.

Zero software licensing costs; low hardware investment amortised 3-5 years;
reusability across course editions; institutional value for TecnoCampus.

\subsection{Removed from Ch.6: Environmental energy calculation detail}

% Reason: Detailed kWh/CO2 calculations with solar breakdown preserved
% in summary form. The comparison table and sustainability bullets are kept.

Development: 150W server + 65W workstation + 15W network = 230W x 820h = 188.6 kWh.
CO2 at Spain avg 0.25 kg/kWh = 47.15 kg. Solar (5kWp, 6634.9 kWh/2025):
34\% offset = 64.1 kWh solar, adjusted CO2 = 31.95 kg (32\% reduction).
Operational: 230W x 360h = 82.8 kWh/year = 20.7 kg CO2/year.

Software costs table (all 0 euro): EVE-NG CE, Parrot OS, Proxmox VE,
Metasploitable, DVWA, VSCode, LaTeX, Git+GitHub (Student Pack), Burp Suite CE,
Wireshark. (Kept in budget tables as a merged row.)

Other costs detail: Electricity 820h x 0.15kWh x 0.20 euro/kWh = 25 euro;
Internet 9 months x 40 euro = 360 euro; Contingency 5\% = 635 euro; Total 1020 euro.

% ============================================================
\section{Chapter 8 --- Development: Lab Detailed Descriptions}
% ============================================================

\subsection{Removed from Ch.8 Labs: Detailed descriptions (replaced by web reference)}

% The following content was condensed to compact lab sections.
% Full content is available in:
%   docs/web/docs/labs/lab1/  (index, infra, walkthrough, vyos, pfsense, server)
%   docs/web/docs/labs/lab2/  (index, infra, walkthrough, synapse app docs)
%   docs/web/docs/labs/lab3/  (index, infra, walkthrough, server-web, server-db)
% Removed from memory 2026-04-27.

\subsubsection{Lab 1 sections removed}
Infrastructure Design Decisions (Layered Network Architecture, DNS-Based Access,
NAT Port Forwarding, CTF-Style Progression), Target Configuration (PEBCAK Corp
service inventory), Student Workflow (Steps 1-5 + Optional Extension),
Persistent Environment Configuration (systemd-networkd, NetworkManager, udev rules),
Results and Conclusions (Lessons Learned narrative).

\subsubsection{Lab 2 sections removed}
Infrastructure Design Decisions (Segmented DMZ, Two-Server Chain, Internal DNS,
OWASP-Aligned Chain), Infrastructure Configuration (VyOS-LAB2, pfSense-LAB2,
Server-A-SYNAPSE, Server-B-SYNAPSE full configs), Vulnerable Applications detail
(SYNAPSE Portal architecture, Monitor API), Attack Walkthrough (Steps 1-7 full
commands), Mitigation and Defensive Analysis (Output Encoding/CSP, Parameterised
Queries, Safe Subprocess, Secret Management), Results and Conclusions,
Implementation Session notes (2026-04-24/25).

\subsubsection{Lab 3 sections removed}
Infrastructure Design Decisions (Role Reversal, Pre-staged Attack, Isolated
Two-Tier, Single Point of Entry), Infrastructure Configuration (VyOS-LAB3,
pfSense-LAB3, Server-Web, Server-DB full configs), Simulated Attack Scenario
narrative, Evidence Walkthrough (Steps 1-7 with full commands and expected output),
Mitigation and Defensive Analysis (SSH Hardening, Sudoers Least Privilege,
File Integrity Monitoring), Results and Conclusions, Lessons Learned.

\end{document}
